기룡이는 KGUPC의 문제 지문을 여러 파일로 나누어 작성하고 있다. 열심히 수정한 내용이 사라지지 않도록 틈틈이 저장하고 있지만, 저장 명령을 실행한다고 항상 파일이 새로 저장되는 것은 아니다.

편집기에는 $1$번부터 $N$번까지의 파일이 열려 있다. 처음에는 모든 파일이 저장된 상태이며, 다음 세 종류의 작업을 수행할 수 있다.

\begin{itemize}
\item \texttt{EDIT i}: $i$번 파일을 수정한다. 이 파일은 저장되지 않은 수정 사항이 있는 상태가 된다.
\item \texttt{SAVE i}: $i$번 파일에 저장되지 않은 수정 사항이 있다면 파일을 저장한다. 그렇지 않다면 아무 일도 일어나지 않는다.
\item \texttt{SAVEALL}: 저장되지 않은 수정 사항이 있는 모든 파일을 각각 저장한다.
\end{itemize}

한 파일을 여러 번 수정했더라도, 한 번 저장하면 그때까지의 수정 사항이 모두 저장된다. 파일을 저장한 뒤 다시 수정하면 새로운 저장이 필요하다. 또한 \texttt{SAVEALL}로 여러 파일을 저장하면 각 파일을 저장할 때마다 저장 횟수가 $1$씩 증가한다.

기룡이가 수행한 작업들이 순서대로 주어질 때, 파일 저장이 수행된 총횟수를 구해 보자. 마지막 작업이 끝난 뒤에는 추가로 저장하지 않는다.